\newpage
\section{\tkzClass{path} 类}

这个类是应 \code{tex.stackexchange.com} 上的一个问题而开发的 [\href{https://tex.stackexchange.com/questions/741663/fill-a-space-made-by-three-arcs/741758#741758}{fill-a-space-made-by-three-arcs}]。

宏包中原本缺少这个概念，不过在某些情况下可以借助 \TIKZ{} 路径来绕过。在创建 \code{conic} 类时，我经常需要操作坐标表，于是把它们形式化为 \code{path} 就变得很自然。

\directlua{
 z.O = point(0, 0)
 z.A = point(8, 0)
 z.E = z.O:rotation(2*math.pi / 3, z.A)
 z.C = z.O:rotation(-2*math.pi / 3, z.A)
 z.F, z.D, z.B = z.O:rotation(math.pi / 3, z.A, z.E, z.C)
 L.AC = line(z.A, z.C)
 L.OB = line(z.O, z.B)
 L.OD = line(z.O, z.D)
 L.CE = line(z.C, z.E)
 L.AE = line(z.A, z.E)
 L.OF = line(z.O, z.F)
 z.G = intersection(L.AC, L.OB)
 z.H = intersection(L.CE, L.OD)
 z.I = intersection(L.AE, L.OF)
 C.GA = circle(z.G, z.A)
 C.FE = circle(z.F, z.E)
 C.BA = circle(z.B, z.A)
 C.IE = circle(z.I, z.E)
 C.DE = circle(z.D, z.E)
 C.FA = circle(z.F, z.A)
 C.HE = circle(z.H, z.E)
 C.GC = circle(z.G, z.C)
 _,z.J = intersection(C.GA, C.FE)
 _,z.K = intersection(C.BA, C.IE)
 _,z.L = intersection(C.DE, C.IE)
 _,z.M = intersection(C.FA, C.HE)
 z.N = intersection(C.BA, C.HE)
 z.P = intersection(C.DE, C.GC)
 % the paths
 local nb = 40
 PA.th1 = C.IE:path(z.E, z.L, nb) + C.DE:path(z.L, z.O, nb) - C.FA:path(z.E, z.O, nb)
 PA.th2 = C.IE:path(z.G, z.A, nb) - C.FA:path(z.M, z.A, nb) - C.HE:path(z.G, z.M, nb)
 PA.th3 = C.IE:path(z.H, z.K, nb) + C.BA:path(z.K, z.C, nb) - C.GC:path(z.H, z.C, nb)
 PA.th4 = C.GA:path(z.A, z.J, nb) + C.FA:path(z.J, z.O, nb) - C.BA:path(z.A, z.O, nb)
 PA.th5 = C.HE:path(z.C, z.N, nb) + C.BA:path(z.N, z.O, nb) - C.DE:path(z.C, z.O, nb)
 PA.th6 = C.HE:path(z.I, z.E, nb) - C.DE:path(z.P, z.E, nb) - C.GA:path(z.I, z.P, nb)}
\begin{center}
  \begin{tikzpicture}[scale = .5, rotate = -30]
    \tkzGetNodes
    \tkzFillCircle[fill=yellow!20](O,A)
    \tkzDrawCircle[ultra thick,black](O,A)
    \tkzClipCircle(O,A)
    \tkzDrawCoordinates[draw,ultra thick,fill=green!15](PA.th1)
    \tkzDrawCoordinates[draw,ultra thick,fill=green!15](PA.th2)
    \tkzDrawCoordinates[draw,ultra thick,fill=blue!15](PA.th3)
    \tkzDrawCoordinates[draw,ultra thick,fill=blue!15](PA.th4)
    \tkzDrawCoordinates[draw,ultra thick,fill=red!15](PA.th5)
    \tkzDrawCoordinates[draw,ultra thick,fill=red!15](PA.th6)
    \tkzDrawCircle[ultra thick,black](O,A)
    \end{tikzpicture}
\end{center}

\begin{verbatim}

\directlua{
 z.O = point(0, 0)
 z.A = point(8, 0)
 z.E = z.O:rotation(2*math.pi / 3, z.A)
 z.C = z.O:rotation(-2*math.pi / 3, z.A)
 z.F, z.D, z.B = z.O:rotation(math.pi / 3, z.A, z.E, z.C)
 L.AC = line(z.A, z.C)
 L.OB = line(z.O, z.B)
 L.OD = line(z.O, z.D)
 L.CE = line(z.C, z.E)
 L.AE = line(z.A, z.E)
 L.OF = line(z.O, z.F)
 z.G = intersection(L.AC, L.OB)
 z.H = intersection(L.CE, L.OD)
 z.I = intersection(L.AE, L.OF)
 C.GA = circle(z.G, z.A)
 C.FE = circle(z.F, z.E)
 C.BA = circle(z.B, z.A)
 C.IE = circle(z.I, z.E)
 C.DE = circle(z.D, z.E)
 C.FA = circle(z.F, z.A)
 C.HE = circle(z.H, z.E)
 C.GC = circle(z.G, z.C)
 _,z.J = intersection(C.GA, C.FE)
 _,z.K = intersection(C.BA, C.IE)
 _,z.L = intersection(C.DE, C.IE)
 _,z.M = intersection(C.FA, C.HE)
 z.N = intersection(C.BA, C.HE)
 z.P = intersection(C.DE, C.GC)
 % the paths
 local nb = 40
 PA.th1 = C.IE:path(z.E, z.L, nb) + C.DE:path(z.L, z.O, nb) - C.FA:path(z.E, z.O, nb)
 PA.th2 = C.IE:path(z.G, z.A, nb) - C.FA:path(z.M, z.A, nb) - C.HE:path(z.G, z.M, nb)
 PA.th3 = C.IE:path(z.H, z.K, nb) + C.BA:path(z.K, z.C, nb) - C.GC:path(z.H, z.C, nb)
 PA.th4 = C.GA:path(z.A, z.J, nb) + C.FA:path(z.J, z.O, nb) - C.BA:path(z.A, z.O, nb)
 PA.th5 = C.HE:path(z.C, z.N, nb) + C.BA:path(z.N, z.O, nb) - C.DE:path(z.C, z.O, nb)
 PA.th6 = C.HE:path(z.I, z.E, nb) - C.DE:path(z.P, z.E, nb) - C.GA:path(z.I, z.P, nb)}
\begin{center}
  \begin{tikzpicture}[scale = .5, rotate = -30]
    \tkzGetNodes
    \tkzFillCircle[fill=yellow!20](O,A)
    \tkzDrawCircle[ultra thick,black](O,A)
    \tkzClipCircle(O,A)
    \tkzDrawCoordinates[draw,ultra thick,fill=green!15](PA.th1)
    \tkzDrawCoordinates[draw,ultra thick,fill=green!15](PA.th2)
    \tkzDrawCoordinates[draw,ultra thick,fill=blue!15](PA.th3)
    \tkzDrawCoordinates[draw,ultra thick,fill=blue!15](PA.th4)
    \tkzDrawCoordinates[draw,ultra thick,fill=red!15](PA.th5)
    \tkzDrawCoordinates[draw,ultra thick,fill=red!15](PA.th6)
    \tkzDrawCircle[ultra thick,black](O,A)
    \end{tikzpicture}
\end{center}

\end{verbatim}

\medskip

除了从一个点列表手动构造路径外，你还可以利用各个几何类自动生成路径。这在绘制圆弧、线段或插值曲线时特别有用。目前，以下类提供了生成路径对象的内置方法：

\begin{itemize}
\item \tkzClass{line} —— 两点之间的线段插值，见 [\ref{ssub:line_path}]
\item \tkzClass{triangle} —— 作为路径的三角形轮廓，见 [\ref{ssub:method_tkzmeth_triangle_path}]
\item \tkzClass{circle} —— 圆周上两点之间的圆弧，见 [\ref{ssub:circle_path}]
\item \tkzClass{conic} —— 参数化的弧或整条曲线，见 [\ref{ssub:conic_path}]
\end{itemize}

这些方法通常返回一个 \tkzClass{path} 对象，其点数为默认或用户指定。例如：

\begin{verbatim}
PA.arc = C.OA:path(z.A, z.B, 40) -- circular arc from A to B on circle OA
PA.seg = L.AB:path(10) -- line segment from A to B with 10 points
PA.tri = T.ABC:path() -- triangle outline path
PA.co = CO.EL:points(0, 1, 50) -- conic arc from t=0 to t=1
\end{verbatim}

这些生成的路径可以合并、反向、填充，或用于 \TIKZ{} 的装饰。


\subsection{概述}

路径对象表示一个有序的二维复点序列，以 \texttt{"(x,y)"} 形式的、与 TikZ 兼容的字符串存放。它支持平移、位似、旋转等几何运算，并可通过 Lua 运算符重载与其他路径组合。

变量 \code{PA} 持有一个表，用于存放路径。它是可选的，你可以自由选择变量名。不过，为清晰和一致起见，推荐使用 \code{PA}。如果你使用自定义变量（例如 Paths），则必须手动初始化它。如果使用了 \code{PA}，\tkzFct{tkz-elements}{init\_elements()} 函数会重新初始化该表。

在处理路径对象时——尤其是涉及若干中间曲线时——强烈建议使用带索引的名字，如 \code{PA.p1}、\code{PA.p2} 等，而不要用一个笼统的名字。

\begin{mybox}
\textbf{规则——优先使用显式的、带索引的名字：}
使用 \code{PA.p1}、\code{PA.p2}、\code{PA.result} 等可提高清晰度、避免意外覆盖，并确保所有几何数据都能通过一次 \tkzFct{tkz-elements}{init\_elements()} 调用重置。
\end{mybox}

这种做法与 \tkzNamePack{tkz-elements} 的整体理念一致：用 \code{PA}、\code{T}、\code{C}、\code{L} 等具名表来存放结构化的几何信息。带索引的命名也有助于让图更易读、可复现、易调试。


它支持用于组合和操作路径的算术运算，以及平移、旋转、位似和反向等几何变换。

\subsection{说明}

\begin{itemize}
\item 点在内部以 \texttt{"(x,y)"} 这样的字符串存放，解析通过工具函数 parse\_point 完成。

\begin{verbatim}
local x, y = utils.parse_point("(3.5, -2)")
-- x = 3.5, y = -2.0
\end{verbatim}


\item 数字格式化（例如为了与 \TIKZ{} 兼容）应由 utils\_checknumber 处理，假定它已在全局或某个工具模块中定义。

\item 该类设计为轻量型，并与 TikZ/LuaLaTeX 工作流兼容。

\end{itemize}





\subsection{构造器 \tkzMeth{path}{path(<list>)}}

创建一个新路径。

\medskip

\texttt{参数： }
 \code{<list>}（可选）：一个 Lua 表，包含以 TikZ 格式字符串给出的点坐标，
        例如 \code{"(0,0)"}、\code{"(1,2)"}。


\begin{center}
  \begin{verbatim}
 PA.name = path(table_of_points) -- Creates a new path object.
  \end{verbatim}
\end{center}

\medskip

若未提供参数，则创建一个空路径。


若提供了数据，输入应是一个以字符串书写的点表，例如 |{ "(0,0)", "(1,0)"|、|"(1,1)" }|，否则创建一个空路径。

下面是一个表示简单三角形的 \tkzname{path}：

  \begin{mybox}
    \begin{verbatim}
      PA.name = path({ "(0, 0)", "(1, 0)","(1, 1)" })
      tex.print(tostring(PA.name))
    \end{verbatim}
  \end{mybox}



\subsection{运算符重载；元方法}

\subsubsection{元方法一览}
\bgroup
  \catcode`_=12
  \small
  \captionof{table}{vector 类的方法。}\label{path:metamethods}
  \begin{tabular}{ll}
  \toprule
   \texttt{元方法} & \texttt{参考} \\
   \midrule
  \tkzMeta{path}{add(path1,path2)}   & [\ref{ssub:metamethod_add}]  \\
  \tkzMeta{path}{sub(path1,path2)}   & [\ref{ssub:metamethod_sub}]  \\
  \tkzMeta{path}{unm(path1)}         &  [\ref{ssub:metamethod_unm}]       \\
  \tkzMeta{vector}{tostring(path1)}   & [\ref{ssub:metamethod_tostring}]  \\
  \bottomrule
  \end{tabular}
\egroup

\subsubsection{元方法 \tkzMeta{path}{add}}
\label{ssub:metamethod_add}
若 p1 和 p2 是两条路径，可以用
  |p3 = p1:add(p2)|
得到第三条路径。

 更方便地，可使用运算符

 \code{p3 = p1 + p2} —— \textbf{拼接} \\
 把 \code{p2} 接在 \code{p1} 之后，返回一条新路径。

\subsubsection{元方法 \tkzMeta{path}{unm}}
\label{ssub:metamethod_sub}
|p2 = p1:unm()|

 更方便地，可使用运算符

  \code{-p} —— \textbf{反向} \\
  返回该路径按相反顺序的一个副本。

\subsubsection{元方法 \tkzMeta{path}{sub}}
\label{ssub:metamethod_unm}
若 p1 和 p2 是两条路径，可以用
|p3 = p1:sub(p2)|
得到第三条路径。

 更方便地，可使用运算符

  \code{p1 - p2} —— \textbf{相减} \\
  等价于 \code{p1 + (-p2)}（把 \code{p1} 与反向的 \code{p2} 拼接）。

\subsubsection{元方法 \tkzMeta{path}{tostring}}
\label{ssub:metamethod_tostring}
字符串表示

\code{tostring} 元方法提供一种可读形式，主要用于调试：

\begin{verbatim}
tostring(p) --> path: { (x1,y1) , (x2,y2) , ... }
\end{verbatim}



\subsection{方法}

\bgroup
\catcode`\_ = 12
\small
\captionof{table}{\texttt{path} 类的方法。}\label{path:methods}
\begin{tabular}{ll}
\toprule
\texttt{方法} & \texttt{参考} \\
\midrule
\tkzMeth{path}{add\_point(z)}     & [\ref{ssub:method_tkzmeth_path_add__point}]\\
\tkzMeth{path}{get(i)}            & [\ref{ssub:method_tkzmeth_path_get}] \\
\tkzMeth{path}{copy()}            & [\ref{ssub:method_tkzmeth_path_copy}] \\
\tkzMeth{path}{count()}           & [\ref{ssub:method_tkzmeth_path_count}] \\
\tkzMeth{path}{translate(dx, dy)} & [\ref{ssub:method_tkzmeth_path_translate}] \\
\tkzMeth{path}{scale(kx, ky)}     & [\ref{ssub:path_scale}] \\
\tkzMeth{path}{homothety(pt, k)}  & [\ref{ssub:path_homothety}] \\
\tkzMeth{path}{map(fx, fy)}        & [\ref{ssub:path_map}] \\
\tkzMeth{path}{rotate(pt, an)}    & [\ref{ssub:method_tkzmeth_path_rotate}] \\
\tkzMeth{path}{close()}           & [\ref{ssub:method_tkzmeth_path_close}]\\
\tkzMeth{path}{sub(i1, i2)}       & [\ref{ssub:method_tkzmeth_path_sub_i1_i2}] \\
\tkzMeth{path}{show()}            & [\ref{ssub:method_tkzmeth_path_show}] \\
\tkzMeth{path}{add\_pair\_to\_path(z1, z2, n)} & 或 add\_pair [\ref{ssub:subsubsection_name}] \\
\tkzMeth{path}{concat(sep)}       &  [\ref{ssub:path_concat}]\\
\bottomrule
\end{tabular}
\egroup



\subsubsection{方法 \tkzMeth{path}{add\_point(pt,<n>)}}
\label{ssub:method_tkzmeth_path_add__point}

该方法向路径追加一个点。

\medskip

\textbf{参数：}
\begin{itemize}
  \item \code{pt}：可以是
  \begin{itemize}
    \item 一个 \code{point} 对象，或
    \item 表示一个点坐标的字符串，例如 \code{"(2,3)"} 或 \code{"(1e-5,0)"}。
  \end{itemize}
  \item \code{n}（可选）：为 TikZ 输出格式化坐标时所用的小数位数。
\end{itemize}

\medskip

若 \code{pt} 是字符串，则用 \code{utils.parse\_point} 解析它。然后该点被格式化（经由 \code{utils.format\_point}）并存入路径。

\medskip

若给出 \code{n}，坐标在存入路径前先按 \code{n} 位小数舍入。若省略，则使用 10 位小数，以便从路径读回的点在进一步计算中保留足够精度。

\medskip

\texttt{注意： }
参数 \code{n} 只影响生成的 TikZ 代码中的数值表示。它不改变内部的几何计算。

\medskip

\textbf{示例}

\begin{tkzexample}[latex=.35\textwidth]
\directlua{
  PA.triangle = path()
  PA.triangle:add_point(point(0,0),0)
  PA.triangle:add_point(point(1,0),0)
  PA.triangle:add_point(point(1,1),0)
  tex.print(tostring(PA.triangle))
}
\end{tkzexample}

\begin{tkzexample}[latex=.35\textwidth]
\directlua{
  PA.p = path({ "(0,0)", "(1,0)", "(1,1)" })
  PA.p:add_point("(2,3)",0)
  PA.p:add_point("(1e-5,0)",5)
  tex.print(tostring(PA.p))
}
\end{tkzexample}
\subsubsection{方法 \tkzMeth{path}{get(i)}}
\label{ssub:method_tkzmeth_path_get}

方法 \code{get(i)} 返回存放在路径 \texttt{PA.p} 中的第 i 个点。

tkz-elements 中的路径把点以格式化的坐标字符串存放（例如 \texttt{"(1.000,2.000)"}），以与 TikZ 兼容并高效地从 Lua 传递到 \TeX。当调用 \code{get(i)} 时，这些坐标会自动转换回一个真正的 Lua 点对象。

\medskip
\texttt{返回值： }

与路径第 i 个元素对应的点对象。

\medskip
\texttt{错误： }

若索引超出有效范围，则引发错误。



\subsubsection{方法 \tkzMeth{path}{copy()}}
\label{ssub:method_tkzmeth_path_copy}

返回该路径的深拷贝。

\begin{tkzexample}[latex=.3\textwidth]
\directlua{
  PA.p1 = path({ "(0, 0)", "(1, 0)","(1, 1)" })
  PA.p2 = PA.p1:copy()
  PA.p2:add_point(point(0, 0),0)}
\begin{tikzpicture}
  \tkzDrawCoordinates(PA.p2)
\end{tikzpicture}
\end{tkzexample}

\subsubsection{方法 \tkzMeth{path}{count()}}
\label{ssub:method_tkzmeth_path_count}

返回路径中点的个数。

\begin{tkzexample}[latex=.3\textwidth]
\directlua{
  PA.p1 = path({ "(0, 0)", "(1, 0)","(1, 1)" })
  PA.p2 = PA.p1:copy()
  PA.p2:add_point(point(0, 0),0)
  tex.print(PA.p2:count())
  }
\end{tkzexample}



\subsubsection{方法 \tkzMeth{path}{translate(dx,dy)}}
\label{ssub:method_tkzmeth_path_translate}

把路径按给定向量 \((dx,dy)\) 平移。

\begin{tkzexample}[latex=.3\textwidth]
\directlua{
  PA.p1 = path({ "(0,0)", "(1,0)", "(1,1)" })
  PA.p2 = PA.p1:translate(2,1)}
\begin{tikzpicture}
\tkzDrawCoordinates(PA.p2)
\end{tikzpicture}
\end{tkzexample}

\subsubsection{方法 \tkzMeth{path}{scale(kx,ky)}}
\label{ssub:path_scale}

该方法相对于原点对路径施加各向异性缩放。

\medskip

\texttt{参数： }
\begin{itemize}
  \item \code{kx}：水平缩放因子，
  \item \code{ky}：竖直缩放因子。
\end{itemize}

\medskip

\texttt{描述： }

路径中的每个点 $P(x,y)$ 变换为
\[
P'(kx \cdot x,\; ky \cdot y).
\]

该变换相对于原点 $(0,0)$ 进行。

\medskip

\texttt{返回： }

一个新的 \code{path} 对象，包含缩放后的点。
原路径保持不变。

\medskip

\texttt{示例： }

\begin{verbatim}
PA.scaled = PA.original:scale(2,0.5)
\end{verbatim}

这把 $x$ 坐标放大一倍、$y$ 坐标缩小一半。

\subsubsection{方法 \tkzMeth{path}{homothety(k,<center>)}}
\label{ssub:path_homothety}

该方法对路径施加均匀缩放（位似）。

\medskip

\textbf{参数}
\begin{itemize}
  \item \code{k}：缩放因子，
  \item \code{center}：可选参考点。
\end{itemize}

\medskip

\textbf{第 5 版的变化}

\code{center} 参数现在是可选的。
若省略，则使用原点 $(0,0)$。

\medskip

\textbf{描述}

每个点 $P(x,y)$ 变换为
\[
P'(x',y') = (x_c + k(x - x_c),\; y_c + k(y - y_c)).
\]

\medskip

\textbf{返回}

一个新的 \code{path} 对象。

\subsubsection{方法 \tkzMeth{path}{map(fx,<fy>)}}
\label{ssub:path_map}

该方法对路径中的每个点施加一个坐标变换。

\medskip

\texttt{参数： }
\begin{itemize}
  \item \code{fx}：作用于 $x$ 坐标的 Lua 函数，
  \item \code{fy}：可选，作用于 $y$ 坐标的 Lua 函数。
\end{itemize}

若省略 \code{fy}，则函数 \code{fx} 同时作用于两个坐标。

\medskip

\texttt{描述： }

路径中的每个点 $P(x,y)$ 变换为
\[
P'(x',y')
\]
满足：
\[
x' = fx(x),
\qquad
y' = fy(y).
\]

若只提供一个函数：
\[
(x,y) \mapsto (f(x), f(y)).
\]

\medskip

\texttt{返回： }

一个新的 \code{path} 对象，包含变换后的点。
原路径不被修改。

\medskip

\texttt{示例： }

对两个坐标施加同一个函数：

\begin{verbatim}
PA.new = PA.old:map(math.abs)
\end{verbatim}

施加不同的变换：

\begin{verbatim}
PA.new = PA.old:map(
  function(x) return 2*x end,
  function(y) return 0.5*y end
)
\end{verbatim}

这产生变换
\[
(x,y) \mapsto (2x, 0.5y).
\]




\subsubsection{方法 \tkzMeth{path}{rotate(pt, an)}}
\label{ssub:method_tkzmeth_path_rotate}

  把路径绕 \(pt\) 旋转角度 \(\theta\)（弧度）。

\begin{tkzexample}[latex=.3\textwidth]
\directlua{
   z.O = point(1, 1)
   PA.base = path({ "(1,0)", "(2,0)", "(2,1)" })
   PA.rotated = PA.base:rotate(z.O, math.pi / 2)}
\begin{tikzpicture}
  \tkzGetNodes
  \tkzDrawCoordinates(PA.base)
  \tkzDrawCoordinates(PA.rotated)
  \tkzDrawPoint(O)
\end{tikzpicture}
\end{tkzexample}

\subsubsection{方法 \tkzMeth{path}{close()}}
\label{ssub:method_tkzmeth_path_close}

在末尾重复第一个点，从而闭合路径。

\begin{tkzexample}[latex=.3\textwidth]
\directlua{
  PA.p = path({ "(0,0)", "(1,0)", "(1,1)" })
  PA.closed = PA.p:close()}
\begin{tikzpicture}
  \tkzDrawCoordinates[fill=green!20](PA.closed)
\end{tikzpicture}
\end{tkzexample}

\subsubsection{方法 \tkzMeth{path}{sub(i1, i2)}}
\label{ssub:method_tkzmeth_path_sub_i1_i2}

返回两个索引之间的一段子路径。


\begin{tkzexample}[latex=.3\textwidth]
\directlua{
  PA.full = path({ "(0,0)", "(1,0)", "(1,1)", "(0,1)", "(0,0)" })
  PA.part = PA.full:sub(2, 4)}
\begin{tikzpicture}
  \tkzDrawCoordinates[fill=green!20](PA.part)
\end{tikzpicture}
\end{tkzexample}


\subsubsection{方法 \tkzMeth{path}{show()}}
\label{ssub:method_tkzmeth_path_show}

在终端打印该路径（用于调试）。

\begin{tkzexample}[latex=.45\textwidth]
\directlua{
  PA.p = path({ "(0,0)", "(1,0)", "(1,1)" })
  PA.p:show()}
\end{tkzexample}

% ADD(AM, 4.25) 补充内容
\subsubsection{方法 \tkzMeth{path}{add\_pair\_to\_path(p, p, n)} 或 \tkzMeth{path}{add\_pair(p, p, n)}}
\label{ssub:subsubsection_name}
\begin{verbatim}
  path:add_pair(z1, z2, decimals)
\end{verbatim}

或等价地：

\begin{verbatim}
 path:add_pair_to_path(z1, z2, decimals)
\end{verbatim}

\medskip
\texttt{用途： }

 取出 \code{z1} 和 \code{z2} 的坐标，按固定小数位数 $n$ 格式化，组合成字符串 \code{"x1/y1/x2/y2"} 追加到路径。然后该字符串被插入当前路径表。该方法向一个 \tkzname{path} 对象追加一个格式化好的点对。它主要用于存放线段，把两个点的坐标成对地放在一起。

\medskip
\texttt{参数： }

\begin{itemize}
  \item \code{z1}、\code{z2} – 以 \code{\{re, im\}} 表给出的点。
  \item \code{decimals}（可选） – 小数位数（默认：\code{10}）。
\end{itemize}

\medskip
\texttt{返回： }

无返回值（就地修改路径）。

\medskip
\texttt{用法示例： }

\begin{verbatim}
   local PA.p = path() -- empty path
   local z.A = point(1, 2)
   local z.B = point(3, 4)
   PA.p:add_pair(z.A, z.B, 2)
   -- PA.p = {"1.00/2.00/3.00/4.00"}
\end{verbatim}


\subsubsection{方法 \tkzMeth{path}{concat(sep)}}
\label{ssub:path_concat}

该方法返回一个字符串，它由一个 \tkzClass{path} 对象中存放的所有坐标拼接而成。这些坐标按顺序用一个分隔符连接。

\medskip
在内部，一个 \tkzClass{path} 包含一个有序的 Lua 点表，以 TikZ 语法书写，例如 \verb|(x,y)|。\tkzMeth{path}{concat} 方法只是把这些条目连接成一条文本路径。

\medskip
\texttt{语法： }
\begin{verbatim}
S = PA.myPath:concat(sep)
\end{verbatim}

\medskip
\texttt{参数： }
\begin{itemize}
  \item \code{sep} – 可选字符串，用于分隔各坐标（默认：单个空格 \verb|" "|）。
\end{itemize}

\medskip
\texttt{返回： }
一个 Lua 字符串，包含该路径的所有点，以 \code{sep} 分隔。所得字符串可直接插入一条 TikZ 路径。

\texttt{示例： }

\begin{tkzexample}[code only]
\directlua{
init_elements()
local A = point(0, 0)
local B = point(2, 1)
local C = point(3, 0)
local P = path:new()
P:add_point(A)
P:add_point(B)
P:add_point(C)
local S = P:concat(" -- ")
}
\end{tkzexample}


\texttt{用途： }

本方法在以下场合很有用：
\begin{itemize}
  \item 从 Lua 动态构造 TikZ 路径；
  \item 导出或组合几何路径；
  \item 生成诸如 \verb|\tkzDrawPath| 之类宏的参数。
\end{itemize}


\subsubsection{示例：准圆}
\label{ssub:director_circle}


圆锥曲线用 \tkzClass{path} 工具绘制。

\begin{minipage}{.55\textwidth}
  \begin{verbatim}
   \directlua{
    init_elements()
    z.O = point(0, 0)
    z.F1 = point(4, 1)
    z.F2 = point(-4, -1)
    z.H = point(4 * math.sqrt(2), 0)
    local a = tkz.length(z.O, z.H)
    CO.EL = conic(EL_bifocal(z.F2, z.F1, a))
    PA.curve = CO.EL:points(0, 1, 50)
    z.A = CO.EL.covertex
    T.HOA = triangle(z.H, z.O, z.A)
    z.P = T.HOA:parallelogram()
    C.OP = circle(z.O, z.P)
    z.L = C.OP:point(0.25)
    T.LJ, T.LK = CO.EL:tangent_from(z.L)
    z.J = T.LJ.pb
    z.K = T.LK.pb}
  \end{verbatim}
\end{minipage}
\begin{minipage}{.45\textwidth}
   \directlua{
    init_elements()
    z.O = point(0, 0)
    z.F1 = point(4, 1)
    z.F2 = point(-4, -1)
    z.H = point(4 * math.sqrt(2), 0)
    local a = tkz.length(z.O, z.H)
    CO.EL = conic(EL_bifocal(z.F2, z.F1, a))
    PA.curve = CO.EL:points(0, 1, 50)
    z.A = CO.EL.covertex
    T.HOA = triangle(z.H, z.O, z.A)
    z.P = T.HOA:parallelogram()
    C.OP = circle(z.O, z.P)
    z.L = C.OP:point(0.25)
    T.LJ, T.LK = CO.EL:tangent_from(z.L)
    z.J = T.LJ.pb
    z.K = T.LK.pb}
\begin{center}
  \begin{tikzpicture}[scale = .4]
   \tkzGetNodes
   \tkzDrawPoints(F1,F2,O)
   \tkzDrawCircles[teal](O,P)
   \tkzDrawPolygon(H,O,A,P)
   \tkzDrawCoordinates[smooth,cyan](PA.curve)
   \tkzDrawSegments[orange](O,P O,L L,J L,K)
   \tkzDrawPoints(F1,F2,O,H,A,P,L,J,K)
   \tkzLabelPoints(F1,F2,O,H,A,P,L,J,K)
  \end{tikzpicture}
\end{center}
\end{minipage}

\begin{verbatim}
\begin{tikzpicture}
 \tkzGetNodes
 \tkzDrawPoints(F1,F2,O)
 \tkzDrawCircles[teal](O,P)
 \tkzDrawPolygon(H,O,A,P)
 \tkzDrawCoordinates[smooth,cyan](PA.curve)
 \tkzDrawSegments[orange](O,P O,L L,J L,K)
 \tkzDrawPoints(F1,F2,O,H,A,P,L,J,K)
 \tkzLabelPoints(F1,F2,O,H,A,P,L,J,K)
\end{tikzpicture}
\end{verbatim}



\subsubsection{经典抛物线}
\label{ssub:classic_parabola}

本示例脱离了欧几里得几何的概念，但所用方法可以复用。这里我们想画一条对称轴平行于 y 轴的抛物线。已知它的三个点，就可以确定方程。绘制抛物线所需的点集，用 \tkzClass{path} 类提供的工具得到。

\begin{minipage}{.55\textwidth}
\begin{verbatim}
\directlua{
 init_elements()
 z.a = point(1, 0)
 z.b = point(3, 2)
 z.c = point(0, 2)
 local A, B, C = parabola(z.a, z.b, z.c)

 function f(t0, t1, n)
 	local PA.tbl = path()
 	for t = t0, t1, (t1 - t0) / n do
 		local y = A * t ^ 2 + B * t + C
    local pt = point(t, y)
    PA.tbl:add_point(pt)
 	end
 	return PA.tbl
 end}
 \end{verbatim}
\end{minipage}
\begin{minipage}{.45\textwidth}
  \directlua{
   init_elements()
   z.a = point(1, 0)
   z.b = point(3, 2)
   z.c = point(0, 2)
   local A, B, C = tkz.parabola(z.a, z.b, z.c)

   function def_curve(t0, t1, n)
   local p = path()
   for t = t0,t1,(t1-t0)/n  do
       y = A*t^2+B*t +C
       pt = point(t, y)
       p:add_point(pt)
    end
    return p
   end

   PA.curve = def_curve(-1,3,100)}

  \begin{center}
    \begin{tikzpicture}
       \tkzGetNodes
       \tkzInit[xmin = -2,xmax=4,ymin =-1,ymax=6]
       \tkzDrawX\tkzDrawY
       \tkzDrawPoints[red,size=4pt](a,b,c)
      \tkzDrawCoordinates[smooth,purple](PA.curve)
    \end{tikzpicture}
  \end{center}
\end{minipage}


\begin{verbatim}
\begin{tikzpicture}
 \tkzGetNodes
 \tkzInit[xmin = -2,xmax=4,ymin =-1,ymax=6]
 \tkzDrawX\tkzDrawY
 \tkzDrawPoints[red,size=4pt](a,b,c)
 \draw[smooth] plot coordinates {%
 \directlua{tex.print(f(-1,3,100))}};
\end{tikzpicture}
\end{verbatim}

\subsection{含多条路径的示例}

\begin{tkzexample}[latex=.45\textwidth]
\directlua{
  z.O = point(0, 0)
  z.A = point(5, 0)
  C.OA = circle(z.O, z.A)
  z.S = C.OA.south
  C.SO = circle(z.S, z.O)
  z.B, z.C = intersection(C.OA, C.SO)
  C.BC = circle(z.B, z.C)
  z.D = intersection(C.OA, C.BC)
  C.CD = circle(z.C, z.D)

  local p1 = C.SO:path(z.C, z.B, 50)
  local p2 = C.BC:path(z.C, z.D, 50)
  local p3 = C.CD:path(z.D, z.B, 50)
  thepath = (-p1) + p2 + p3 }
\begin{tikzpicture}[scale=.5]
  \tkzGetNodes
  \tkzDrawCircles(O,A S,O)
  \tkzDrawArc(B,C)(D)
  \tkzDrawArc(C,D)(B)
  \tkzDrawCoordinates[fill = purple!20,
                      opacity=.4](thepath)
  \tkzDrawCoordinates[smooth,red,thick](thepath)
  \tkzDrawPoints(A,O,B,C,S,D)
  \tkzLabelPoints(A,B,C)
\end{tikzpicture}
\end{tkzexample}
\endinput